\chapter{مقدمه}
\newpage
‎\renewcommand{\baselinestretch}{2}\normalsize‎
\label{introduction}
\section{تعریف مساله}



امروزه شرکت‌های زیادی نظیر \lr{IBM} و \lr{Google} رقابت گسترده‌ای در زمینه ساخت کامپیوترهای کوانتومی دارند \cite{cuomo2020towards,cacciapuoti2020entanglement,caleffi2018quantum,cacciapuoti2019quantum}. در طول ٥٠ سال گذشته، روند شگفت‌انگيز کوچک‌سازي\LTRfootnote{\lr{Miniaturization}} درصنعت کامپيوتر پديدار شده است. در حالي که يک ريزپردازنده در سال ١٩٧١ تقريباً شامل ٢٣٠٠ ترانزيستور بود، امروزه ريزپردازنده‌اي در همان اندازه، در برگيرنده‌ی بيش از صدها ميليون، ترانزيستور است. در سال‌های اخیر ابعاد اجزاي مدارهاي مجتمع به حدود چند اتم رسیده است. بدين‌منظور، دانشمندان نظريه جديدي تحت عنوان محاسبات کوانتومي بيان نمودند که به رفتار ذرات در لايه‌هاي اتمي و زير اتمي مي‌پردازد و مجموعه اي از سه علم نظريه اطلاعات،
علوم كامپيوتر و فيزيک کوانتوم است.
بر طبق تعريفي که از محاسبات کوانتومي در منابع مختلف آمده است، اطلاعات و محاسبات کوانتومي مطالعه‌ي عمليات پردازش اطلاعات مي‌باشد که مي‌توانند با استفاده از سيستم‌هاي مکانيک کوانتومي صورت گيرند. مکانيک کوانتومي چارچوبی رياضي با مجموعه‌اي از قوانين براي ساخت تئوري‌هاي فيزيکي مي‌باشد\cite{mcmahon2007quantum}.

رایانه کوانتومی ماشینی است که از پدیده‌ها و قوانین مکانیک کوانتوم مانند برهم‌نهی\LTRfootnote{\lr{Superposition}} و درهم‌تنیدگی\LTRfootnote{\lr{Entanglement}} برای انجام محاسباتش استفاده می‌کند. این رایانه‌ها با رایانه‌های فعلی که با ترانزیستورها کار می‌کنند، تفاوت اساسی دارند. ایده اصلی که در پس رایانه‌های کوانتومی نهفته است این است که می‌توان از خواص و قوانین فیزیک کوانتوم برای ذخیره‌سازی و انجام عملیات روی داده‌ها استفاده نماید.

امروزه محاسبات کوانتومی به دلیل توانایی در افزایش سرعت الگوریتم‌ها توجه زیادی را به خود جلب نموده است. همچنین با پیشرفت روز افزون تکنولوژی، امروزه این محاسبات با محدودیت‌هايي مانند اتلاف گرما و کاهش سرعت مواجه هستند. از طرفی قدرت محاسبات کوانتومی با افزایش تعداد کیوبیت‌های آن به صورت نمایی رشد می‌کند\cite{cacciapuoti2019quantum}. یکی از چالش‌های اساسی در محاسبات کوانتومی این است که کیوبیت‌ها با کوچکترین اغتشاشی در محیط خراب شده و از بین می‌روند، مگر اینکه در دمای بسیار پایین نگه‌داری شوند. یکی از اغتشاشاتی که می‌توان منجر به از بین رفتن اطلاعات کوانتومی و ایجاد ناهمدوسی\LTRfootnote{\lr{Decoherence}} کیوبیت‌ها شود، تعامل آن‌ها در محیط است که با افزایش تعداد آن‌ها، این اطلاعات شکننده و در معرض خطا قرار می‌گیرند. بنابراین ساخت کامپیوترهای کوانتومی در مقیاس بزرگ و با تعداد کیوبیت‌های زیاد بسیار چالش برانگیز، و در عین‌حال دستیابی به آن، هدفی بسیار ارزشمند است.

محاسبات کوانتومي\cite{Nielsen1998} داراي مزايا و روش‌هايي است که باعث مي‌شود مسائلي که حل آن‌ها در محاسبات کلاسیک بسيار سخت است، با الگوریتم‌ها و روش‌هاي خاص محاسبات کوانتومي به آساني و با سرعت بالا حل شوند. کلاس مسائلي که با الگوريتم‌هاي کوانتومي دارای پیچیدگی زمانی چند جمله‌اي مي‌باشند، بزرگتر از کلاس مسائلي است که با يک ماشين تورينگ قطعي يا به عبارتي يک کامپيوتر کلاسیک قابل حل مي‌باشد. به عنوان مثال الگوريتم‌هاي تجزیه\cite{shor-1997-26}، جستجو در يک پايگاه داده‌ي نامرتب\cite{Groverd} و شبيه‌سازي سيستم‌هاي محاسبات کوانتومي از جمله‌ي اين مسائل مي‌باشند. الگوريتم کوانتومي شور \LTRfootnote{\lr{Shor}} جهت تجزیه‌ی اعداد بزرگ يکي از مثال‌هايي است که در پياده‌سازي آن، کامپيوترهاي کوانتومي حتي از بزرگترين ابرکامپيوترهاي کلاسیک هم بسيار سریع‌تر عمل می‌کنند. همچنين محاسبات کوانتومي و اطلاعات کوانتومي در توسعه ابزارهايي براي بهتر کردن درک ما درباره مکانيک کوانتومي و شفاف کردن پيش‌بيني‌هاي آن در ذهن انسان کمک بسياري مي‌کنند.

درحالي که يک کامپيوتر کلاسیک مي‌تواند يک کامپيوتر کوانتومي را شبيه‌سازي کند، اما به روش کارا آن را انجام نمي‌دهد. کامپيوترهاي کوانتومي افزايش سرعت اساسي نسبت به کامپيوترهاي کلاسیک دارند. اين مزيت سرعت، چنان چشمگير است که محققان اعتقاد دارند که هيچ پيشرفت قابل توجهي در کامپيوترهاي کلاسیک، فضاي خالي بين توان يک کامپيوتر کلاسیک و توان يک کامپيوتر کوانتومي را پر نخواهد کرد.

علی‌رغم تمامی مزایایی که کامپیوترهای کوانتومی دارا می‌باشند، اما موانع بزرگی وجود دارند. محققان بر این باورند که قبل از دستیابی به کامپیوتر‌های کوانتومی بزرگ- مقیاس کاربردی، بایستی این مسائل و مشکلات برطرف شود. برای نمونه در فاز محاسباتی عملیات کوانتومی، کوچکترین اغتشاش در سیستم کوانتومی (مثلا یک فوتون سرگردان یا موجی از تشعشع‌های الکترومغناطیسی) باعث قطع شدن محاسبات کوانتومی می‌شود که ناهمدوسی نام دارد. 

یکی از کاربرد‌های جدید امروزی و آینده در محاسبات کوانتومی که مورد توجه طرفداران تکنولوژی بوده و  در آن با بهره ‌برداری از خصوصیات کوانتومی عجیب فوتون‌ها و الکترون‌ها، پیام‌ها و اطلاعات به طور امن ارسال می‌شوند، اینترنت کوانتومی است. این خصوصیات سبب می‌گردند تا دولت‌ها، ارتش‌ها، بانک‌ها و موسسات مالی برای امن نمودن تمام امور خود از جمله قراردادها و تراکنش مالی به اینترنت کوانتومی تمایل نشان دهند. اینترنت کوانتومی یکی از چارچوب‌های اصلی و اساسی در سیستم‌های کوانتومی توزیع شده است که در آن از خاصیت در هم‌تنیدگی کوانتومی استفاده می‌کند\cite{cuomo2020towards}. در این پدیده دو ذره کوانتومی (حتی در فواصل بسیار دور) اطلاعات خود را با یکدیگر به اشتراک می‌گذارند. در واقع با استفاده از جفت فوتون‌های ایجاد شده، می توان اطلاعات کوانتومی را با استفاده از این خاصیت بین دو نقطه جابجا کرد. از این خاصیت در شبکه‌های کوانتومی نظیر اینترنت کوانتومی و در حالت کلی‌تر در سیستم‌های کوانتومی توزیع شده به‌طور فراوان استفاده می‌گردد.

همان‌طور که بالاتر بیان گردید، فعل و انفعال بین کیوبیت‌ها در محیط منجر به ناهمدوسی آن‌ها شده و با افزایش تعداد آن‌ها این اطلاعات شکننده و با خطا مواجه می‌شوند\cite{krojanski2004scaling}. برای غلبه بر مشکل ذکر شده در بالا، می‌توان سیستم کوانتومی را به تعدادی زیر سیستم با تعداد کیوبیت کمتر تجزیه نمود و یک سیستم کوانتومی توزیع شده بوجود آورد. بنابراین ساخت کامپیوتر‌های کوانتومی توزیع شده ضروری است\cite{van2010}.

بدین روی محاسبات کوانتومی توزیع شده راه حل مناسبی برای ساخت سیستم‌های کوانتومی در مقیاس بزرگ است. برای داشتن یک سیستم کوانتومی توزیع شده، نیاز به مکانیزمی برای ایجاد ارتباط بین زیر‌سیستم‌ها است. یکی از پروتکل‌های اصلی در محاسبات کوانتومی توزیع شده، پروتکل دورنوردی\LTRfootnote{\lr{Teleportation}} است که از خاصیت در‌هم‌تنیدگی کیوبیت‌ها برای توزیع نمودن اطلاعات کوانتومی استفاده می‌کند. این پروتکل دارای هزینه‌ی زیادی برای سیستم کوانتومی است. بر طبق تئوری تکثیرناپذیری\LTRfootnote{\lr{No-Cloning}} کوانتومی، زمانی که یک کیوبیت به مقصد خود دورنورد می‌شود، نمی‌تواند در زیر سیستم خود استفاده گردد و تا زمان برگشت کیوبیت دورنورد شده، امکان استفاده از کیوبیت در زیرسیستم خود وجود ندارد. به‌همین دلیل تعداد دورنوردی‌ها در یک سیستم کوانتومی توزیع شده چالش اساسی محسوب می‌گردد. در نتیجه ارائه‌ی یک الگوریتم جهت کاهش این تعداد، کمک قابل توجهی به طراحی و توسعه‌ی دیگر مدارهای کوانتومی خواهد نمود.

شبکه‌های کوانتومی توزیع شده شامل تعدادی سیستم‌های کوانتومی هستند که در فواصل دور از هم قرار دارند و از قوانین و پدیده‌های اساسی مکانیک کوانتومی نظیر برهم‌نهی، در هم‌تنیدگی و اندازه‌گیری کوانتومی جهت انتقال اطلاعات خود استفاده می‌کنند و در کابردهای جدید بسیاری قابل استفاده هستند\cite{loncar2019development}. 
بدین روی جهت داشتن یک کامپیوتر کوانتومی در مقیاس بزرگ، می‌توان شبکه‌ای از کامپیوترهای کوانتومی با ظرفیت محدود (تعداد کیوبیت‌های محدود و مشخص) داشت به‌طوری که ارتباط بین آن‌ها از طریق یک کانال کوانتومی یا کلاسیک و یا هر دو باشد و بتواند رفتار کل سیستم را پیاده‌سازی نماید. این ساختار بعنوان کامپیوتر کوانتومی توزیع شده\LTRfootnote{\lr{Distributed Quantum Circuit}} شناخته می‌شود
\cite{nickerson2013topological}.

 همان‌طور که در تکنولوژی‌های کلاسیک اطلاعات بایستی با یکدیگر در تعامل باشند، در تکنولوژی کوانتومی نیز، نیاز به تبادل اطلاعات و حالت‌های کوانتومی است. برای انجام محاسبات در مدارهای کوانتومی توزیعی، بایستی حالت‌هاي کوانتومی بین دو زیر سیستم کوانتومی جابه‌جا شده و برای انجام این ارتباط بایستی از پروتکل دورنوردی کوانتومی استفاده شود\cite{Bennett}. این پروتکل، پروتکلی اصلی و مهم در سیستم‌های توزیعی بوده و دارای هزینه بالایی است و جهت انتقال حالت یک ذره به ذره‌ای دیگر، از خاصیت درهم‌تنیدگی کوانتومی استفاده می‌کند\cite{Whit} تا بتواند اطلاعات کوانتومی را انتقال دهد. این پروتکل در تکنولوژی‌های بسیاری مانند علم فوتون\LTRfootnote{\lr{Photonics}} \cite{bouwmeester1997}، \lr{NMR} \cite{Nielsen1998} و یون‌های به دام افتاده\cite{Riebe2004}  \LTRfootnote{\lr{Trapped ions}} استفاده می‌گردد. در مدار کوانتومی توزیعی، هر زیرسیستم دارای تعداد مشخص کیوبیت می‌باشد. سیستم توزیع شده از تعدادی دروازه‌های محلی\LTRfootnote{\lr{Local gate}} و سراسری\LTRfootnote{\lr{Global gate}} تشکیل می‌شود. دروازه‌های محلی بر روی کیوبیت‌های هر زیرسیستم اعمال می‌شوند. امّا دروازه‌های سراسری بین دو زیرسیستم قرار گرفته‌اند و نیازمند استفاده از دو کیوبیت در دو زیر سیستم مختلف می‌باشند. 


حال باتوجه به اینکه رفتن کیوبیت به زیر سیستم مقصد و برگشت مجدد آن به زیرسیستم مبدا جهت اجرای دروازه‌های سراسری هزینه‌بر است \cite{sangouard2011quantum}، \cite{duan2001long}
،\cite{g2021efficient}، بایستی بتوان تعداد دورنوردی‌های مورد استفاده برای ساخت مدار توزیعی کوانتومی را به حداقل خود کاهش داد. از طرفی از آنجایی که هدف اصلی ما از طراحی سیستم‌های کوانتومی توزیع شده، داشتن سیستم‌های کوانتومی با تعداد کیوبیت کمتر بوده، بنابراین یکی دیگر از اهداف مهم در طراحی این سیستم‌ها توزیع متوازن تعداد کیوبیت‌ها در زیر سیستم‌ها می‌باشد. %همچنین زمان اجرای این دروازه‌ها اینکه کدام کیوبیت آن‌ها (در دروازه‌های دو کیوبیتی، کیوبیت هدف یا کیوبیت کنترل آن‌ها) دوربر شود، امری بسیار مهم محسوب می‌شود. 

\section{چالش‌های موجود}
در این بخش به بررسی چالش‌های موجود در مساله پرداخته شده است.
\begin{itemize}
\item
همان‌طور که در بالا اشاره گردید، قدرت محاسبات کوانتومی با افزایش تعداد کیوبیت‌های آن افزایش می‌یابد. از طرفی یکی از چالش‌های اساسی در محاسبات کوانتومی تعامل کیوبیت‌ها در محیط است که با افزایش تعداد آن‌ها این اطلاعات شکننده و در معرض خطا قرار می‌گیرند. همچنین با پیشرفت روز افزون تکنولوژی، امروزه این محاسبات با محدودیت‌هایی مانند اتلاف گرما و کاهش سرعت مواجه هستند. بنابراين داشتن یک سیستم کوانتومی در مقیاس بزرگ کاری دشوار است. بدین‌روی محاسبات کوانتومی توزیع شده راه حل مناسبی برای ساخت سیستم‌های کوانتومی در مقیاس بزرگ می‌باشند.
\item
یکی از مهمترین چالش‌ها در محاسبات کوانتومی توزیع شده، عدم وجود یک مدل مناسب جهت ارائه‌ی مدار می‌باشد. در تمامی مطالعات پیشین، هیچ مدلی جهت نمایش محاسبات توزیع شده ارائه نشده بود. در نتیجه ارائه‌ی مدلی خاص جهت انجام و پردازش‌ سیستم‌های توزیع شده احساس می‌شد. از طرفی از آنجایی‌که نیاز است تا کیوبیت‌های سیستم در تعدادی افراز تقسیم شود، نیاز به ارائه الگوریتمی جهت افرازبندی گراف است. بنابراین ارائه‌ی مدل مناسب جهت افرازبندی کیوبیت‌ها با در نظر گرفتن محدودیت‌هاي مدارهای کوانتومی از چالش‌های اساسی محسوب می‌گردد.
 \item
برای داشتن یک سیستم کوانتومی توزیع شده، نیاز به استفاده از پروتکل دورنوردی جهت ایجاد ارتباط بین زیر سیستم‌ها است. این پروتکل از خاصیت درهم‌تنیدگی کیوبیت‌ها برای توزیع نمودن اطلاعات کوانتومی استفاده نموده و دارای هزینه زیادی برای سیستم کوانتومی می‌باشد. همانطور که بالاتر بیان شد، بر طبق تئوری تکثیرناپذیری کوانتومی، زمانی که یک کیوبیت به مقصد خود دورنوردی می‌شود، نمی‌تواند در زیرسیستم خود استفاده گردد. بنابراین تعداد دورنوردی‌ها در یک سیستم کوانتومی توزیع شده چالش اساسی محسوب می‌گردد. در نتیجه، ارائه‌ی یک الگوریتم جهت کاهش این تعداد، کمک قابل توجهی به طراحی و توسعه دیگر مدارهای کوانتومی خواهد نمود.

\item
گاهی اوقات کاهش بیش از اندازه تعداد پروتکل دورنوردی، منجر به عدم وجود توازن در تعداد کیوبیت‌های هر زیر سیستم می‌گردد. بنابراین همزمان با کاهش تعداد دورنوردی‌های کوانتومی توزیع متوازن کیوبیت‌ها نیز هدف اصلی قرار می‌گیرد. بنابراین در طراحی یک سیستم کوانتومی توزیع شده بایستی بتوان تعادلي بين اين دو هدف ايجاد گردد.

\item
از دیگر چالش‌های موجود در مدارهای کوانتومی توزیع شده، این است که پس از توزیع کیوبیت‌ها در تعدادی افراز، در سطح اجرا تصمیم گرفته شود که کیوبیت کنترلی به سمت افرازی که کیوبیت هدف قرار دارد، دورنوردی شود یا بالعکس. 

\item
یکی دیگر از چالش‌ها این است زمانی‌که یک کیوبیت از یک افراز (مبدا) به افراز دیگر (مقصد) دورنوردی می‌شود، چه زمانی کیوبیت دورنوردی شده به مبدا اولیه خود بازگردد. 
\end{itemize}


\section{اهداف مساله}

با توجه به چالش‌های بیان شده در بخش قبل، در این رساله سعی شده است بتوان روشی ارائه داد تا يك مدار یکپارچه کوانتومی را در تعدادي زيرسيستم كوانتومي توزيع نمايد، بطوری که اهداف زیر مقرر گردد.
\begin{itemize}
\item
با توجه به چالش‌های موجود، نبود مدلی مناسب، جهت ارائه محاسبات کوانتومی توزیع شده به شدت احساس می‌شد. در این رساله یکی از اهداف اصلی ارائه مدلی دقیق برای انجام محاسبات توزیع شده است. همچنين كاهش تعداد دورنوردي‌هاي كوانتومي در دوسطح انجام مي‌گردد. در سطح اول سعی شده است كيوبيت‌هاي مدار در تعدادي افراز توزيع شود. براي اين كار، اين مساله را با استفاده از برنامه ریزی خطی مدل نموده و دو الگوريتم جهت افرازبندي آن‌ها در دو افراز و چند افراز ارائه خواهد شد. در سطح دوم به دلیل نیاز به اطلاعات دقیق‌تر مدار، مدلی ماتریسی ارائه گردیده است تا كه در آن اطلاعات اولويتي دروازه‌ها و نوع كنترل يا هدف بودن كيوبيت‌هاي آن‌ها نگهداري شده است. سپس در سطح دوم بهینه‌سازی سعی گردیده است با ارائه تعدادی قوانین ابتکاری سیستم کوانتومی توزیع شده بهینه‌ گردد.
\item
توزیع کیوبیت‌ها در افزارها بایستی متعادل باشد. برای ایجاد این توازن، فاکتور \lr{$\omega$} به نام فاکتور تحمل توازن بار\LTRfootnote{\lr{Load Balance Tolerance}} در نظر گرفته شده است. خروجی الگوریتم سطح اول، برچسب‌گذاري کیوبیت‌ها به صورت\lr{$f:Q\rightarrow \{1,2...,K\}$} مي‌باشد، به‌طوری که نامساوی \eqref{eq15} برقرار باشد.

\begin{equation}
\begin{split}
|\{q\in Q | f(q)=i\}|<=(1+\omega)\frac{|Q|}{K}  \\
\forall i=\{1,2,...,K\}
\end{split}
\label{eq15}
\end{equation}

در این نامعادله، $|Q|$ برابر تعداد کل کیوبیت‌ها، $K$ تعداد افرازها و $f(q)$ بیانگر تابعی است که کیوبیت‌ها را در بین زیر سیستم‌ها افراز می‌کند.
\item
زمانی‌که یک کیوبیت از مبدا خود به مقصد مورد نظر (از یک افراز به افرازی دیگر) دورنوردی می‌شود، چنان‌چه در مبدا اولیه خود مورد نیاز باشد، چه زمانی به مبدا خود بازگردد؟ زمان بازگشت آن به مبدا اولیه امری مهم محسوب می‌گردد. از آنجایی‌که ممکن است با قرار داشتن اين كيوبيت در مقصد، بتوان دروازه‌های بیشتری که به کیوبیت دورنوردی شده نیاز داشته، اجرا شوند. چه بسا که این کیوبیت با بازگشت به مبدا خود بیکار بماند. 
\item
پس از افرازبندی کیوبیت‌ها در سطح اول، چنانچه یک دروازه سراسری باشد، تصمیم‌گیری در مورد اینکه کیوبیت هدف آن به سمت کیوبیت کنترلی دورنوردی شود یا بالعکس، یک چالش اساسی است. این هدف و هدف ذکر شده‌ی قبلی در بهینه سازی سطح دوم مورد استفاده قرار می‌گیرند.

\end{itemize}

\section{جنبه‌ی جدید بودن و نوآوری }

جنبه های جدید بودن و نوآوری سیستم پیشنهادی بشرح زیر است:
\begin{itemize}
\item
بهینه‌سازی تعداد دورنوردی‌ها در دو سطح: در این رساله، بهینه‌سازی در دو سطح انجام می‌گردد. در سطح اول سعی شده است تا بتوان مدار کوانتومی تبدیل به گراف شده و سپس با استفاده از یک الگوریتم بهینه‌سازی، گراف را به $K$ افراز تجزیه نمود. سپس در سطح دوم، با ارائه‌ی تعدادی روش ابتکاری سعی در کاهش تعداد دورنوردی‌ها گردیده است. تنها در یک پژوهش که در کارهای پیشین به آن اشاره خواهد شد\cite{Pablo}، مدار کوانتومی به ابرگراف تبدیل شده و با استفاده از الگوریتم تجزیه ابرگراف، مدار به $K$ افراز تجزیه شده است و هیچ بهینه‌سازی در سطح دوم نداشته است.
\item
در تمامی مطالعات پیشین، عدم وجود مدلی مناسب برای مدارت کوانتومی احساس شده است. در این رساله در سطح اول بهینه‌سازی، مدل برنامه ریز خطی مساله به طور دقیق بیان شده است. و در سطح دوم  یک مدل ماتریسی ارائه شده تا با استفاده از آن بتوان بهینه سازی را انجام داد.

\item
هنگام اجرای یک دروازه‌ی سراسری دو کیوبیتی، تصمیم‌گیری در مورد دورنوردی نمودن کیوبیت هدف به سمت کیوبیت کنترلی و یا بالعکس، تنها در پژوهش \cite{Zomorodi} بررسی شده است. آن‌ها با در نظر گرفتن تمام حالت‌های ممکن بر روی دو افراز تصمیم‌گیری را انجام داده‌اند. اما در این رساله این تصمیم گیری برای چندین افراز و بدون در نظر گرفتن تمامی حالت‌های ممکن انجام می‌پذیرد.
\end{itemize}

\section{ساختار پایان نامه}
اين پايان‌نامه شامل شش فصل است. در فصل اول، كليات، اهداف و ضرورت انجام تحقيق و ايده‌هاي كلي براي مقابله چالش‌ها بيان گردید. در فصل دوم، به بیان مقدمات و آشنایی با مفاهیم سیستم‌های کوانتومی و محاسبات کوانتومی توزیع شده پرداخته می‌شود. فصل سوم کارهاي مختلف انجام شده از نقطه نظرهاي مختلف بررسی شده است. در فصل چهارم، به معرفی راهكارهای پيشنهادي و در فصل پنجم به ارزيابي نتايج بدست آمده پرداخته می‌شود. در نهایت در فصل آخر، نتيجه‌گيري و مسيرهاي اصلي كارهاي آينده بيان خواهد شد.
